\section{Introduction}
\label{sec:introduction}

Information enters financial markets through a heterogeneous network of
public news services, professional terminals, broker desks, and, in
the last decade, real-time social and chat-based aggregators that an
increasing share of retail and small-shop traders monitor as their
primary news feed. These newer channels carry the same headlines that
appear on Bloomberg or Reuters but at a variable latency and with a
publication timestamp that is not the timestamp of the underlying
primary source. Whether and how prices react to the arrival of news
on such channels---and whether the directional content of those
headlines can be extracted in a useful way by modern language
models---is a question with direct practical relevance to a large
and growing audience of intraday traders, and a methodological
question for the event-study literature that has historically focused
either on scheduled macro releases \citep{andersen2003micro} or on
slower-moving newspaper text
\citep{tetlock2007giving, loughran2011liability}.

This paper studies the relation between unscheduled news headlines
from a representative real-time news service---the FinancialJuice
Discord newsfeed---and intraday price reactions in the EUR/USD spot
exchange rate and the Nasdaq-100 index, over a thirteen-month sample
from 24 March 2025 to 5 May 2026. The dataset comprises 2{,}449
hand-filtered ``gold'' news events drawn from 63{,}016 raw messages
and aligned to one-minute price bars on both assets. Each event is
independently labelled by a large language model with a discrete
sentiment for USD and for NDX, an expected magnitude, a surprise
level, and a self-reported confidence. We test fourteen pre-specified
hypotheses covering volatility, direction, persistence, time-of-day,
cross-asset behaviour, and pre-event drift, and we report the central
findings with Benjamini-Hochberg false discovery rate adjustment
\citep{benjamini1995controlling} applied uniformly across the full
hypothesis battery.

Three findings warrant emphasis. First, the news events are robustly
associated with elevated intraday magnitude: the mean absolute return,
the high-low range, and the maximum absolute intra-window move are
between $1.3\times$ and $1.9\times$ their matched-baseline counterparts
across all assets and windows tested, and the result is robust to
winsorisation at the $1\%$ tails, to the introduction of multivariate
controls for category, surprise, magnitude, confidence, cluster size,
and pre-event drift, and to a pre/post split at the LLM's training
data cutoff. The intra-window range and maximum-absolute-move
statistics are particularly robust, with $q$-values below
$10^{-30}$ on every cell. Second, the directional content of LLM-derived sentiment
is modest: direction hit rates range from $49\%$ to $54\%$ on the
primary windows, with only a small number of cells surviving the FDR
correction. A backtest with realistic spread and slippage costs
confirms that this edge is too small to be economically exploitable
in the naive forms tested. Third, the LLM sentiment nonetheless
out-performs both the standard Loughran-McDonald finance dictionary
\citep{loughran2011liability} and the FinBERT-tone neural classifier
\citep{yang2020finbert} on direction hit rate by $2$--$6$ percentage
points across the primary windows, and produces cross-asset
disagreement (different sentiment for USD and for NDX on the same
event) on $73.5\%$ of events---a property that neither single-polarity
baseline can express by construction.

We also report two methodologically substantive findings. The
absolute return in the fifteen minutes preceding a Discord headline
is elevated above the matched baseline by a similar magnitude as in
the post-event window. Without primary-source timestamps we cannot
distinguish whether this reflects information arrival latency in the
public feed or genuine pre-disclosure activity, but the practical
implication for users of the feed is the same: entering a trade at
the timestamp of the public headline is, on average, already late.
And the sign of the $+15$-minute return agrees with the sign of the
$+4$-hour return on $57.6\%$ of events, with median absolute moves
$3$--$4.5\times$ larger over the longer window, so the initial
reaction tends to be extended rather than reversed.

Our contribution is fourfold. We assemble a sizeable intraday
event-study dataset on an actively-used public news feed and
document its construction in enough detail to support replication.
We propose an event-study workflow with matched-baseline
standardisation, explicit cluster dependence handling, and uniform
FDR adjustment as applied to a large battery of hypotheses. We
compare an LLM-based sentiment classifier to two reference
baselines---the Loughran-McDonald finance dictionary and the
FinBERT-tone neural classifier---on identical event windows,
providing direct quantitative support for the LLM methodology
against both classical and BERT-era baselines. And we make the
entire pipeline---including the validation harness, the LLM cache,
the dictionary baseline, the FinBERT baseline, and the three-way
comparison---public in a single repository with seeded reproducibility
(see Reproducibility Statement).

The remainder of the paper is organised as follows.
Section~\ref{sec:literature} reviews the relevant literature on event
studies, financial sentiment analysis, and online news.
Section~\ref{sec:data} describes the data.
Section~\ref{sec:methodology} develops the methodology.
Section~\ref{sec:results} presents the results, grouped by theme.
Section~\ref{sec:discussion} discusses the implications.
Section~\ref{sec:limitations} details the limitations of the design
and the directions for future work, and Section~\ref{sec:conclusion}
concludes.
